% GENERATED FILE. Edit canonical lessons, then run npm run generate:books.
% canonical-source-hash: fnv1a64:a30aa6d2af9f78eb
% canonical-lessons: RU-C132-to-get-dressed, RU-C132-to-wash, RU-C132-to-get-up, RU-C132-to-sit-down, RU-C132-to-fall

\chapter{To get dressed, To wash, To get up, To sit down, To fall}
\label{ch:ru-to-get-dressed-to-wash-to-get-up-to-sit-down-to-fall}
\hlchaptermodality{132}

\begin{chapteropening}
\textbf{By the end of this chapter:} \emph{I can say to get dressed, to wash, to get up, to sit down and to fall.}

Complete the last lesson of chapter 132: I can say to get dressed, to wash, to get up, to sit down and to fall.
\end{chapteropening}

\section[\texorpdfstring{\ru{одеваться} (odevátsya)}{одеваться (odevátsya)}]{\ru{одеваться} (odevátsya) — to get dressed}

\label{lesson:RU-C132-to-get-dressed}



\begin{quote}
Before the new one: say the Russian for to believe, then the Russian for to drive.
\end{quote}

\subsection*{You'll want to know: \ru{одеваться}}
\textbf{\ru{одеваться}} (\emph{odevátsya}) — \textquotedblleft{}to get dressed\textquotedblright{}.

\begin{cousinweb}
\textbf{\ru{одеваться}} is \emph{\ru{о}-}, \textquotedblleft{}around\textquotedblright{}, with \emph{\ru{девать}}, \textquotedblleft{}to put\textquotedblright{}, and \emph{-\ru{ся}}, \textquotedblleft{}oneself\textquotedblright{}: to put clothes around yourself. The root \emph{\ru{де}-}, \textquotedblleft{}to put\textquotedblright{}, is in \emph{\ru{одеяло}}, \textquotedblleft{}blanket\textquotedblright{}.
\end{cousinweb}

\begin{grammarlens}[title={What you've built}]
One more word for this chapter.
\end{grammarlens}

\subsection*{Your turn}
\begin{itemize}
\raggedright
  \item \emph{Say it:} \emph{odevátsya}
  \item \emph{Say it:} \emph{odevátsya}, once more
  \item \emph{Say it:} say \emph{odevátsya}
  \item \emph{Recall:} say the Russian for to believe, then the Russian for to drive, then say \emph{odevátsya} again
\end{itemize}

\begin{tcolorbox}[breakable,colback=teal!4,colframe=teal!35!black,title={Before you move on}]
What does \ru{одеваться} mean? (\textquotedblleft{}To get dressed\textquotedblright{}.) Say it once more.
\end{tcolorbox}
\section[\texorpdfstring{\ru{умываться} (umyvátsya)}{умываться (umyvátsya)}]{\ru{умываться} (umyvátsya) — to wash (oneself)}

\label{lesson:RU-C132-to-wash}



\begin{quote}
Before the new one: say the Russian for to drive, then the Russian for to get dressed.
\end{quote}

\subsection*{You'll want to know: \ru{умываться}}
\textbf{\ru{умываться}} (\emph{umyvátsya}) — \textquotedblleft{}to wash (oneself)\textquotedblright{}.

\begin{cousinweb}
\textbf{\ru{умываться}} is \emph{\ru{у}-} with \emph{\ru{мыть}}, \textquotedblleft{}to wash\textquotedblright{}, and \emph{-\ru{ся}}, \textquotedblleft{}oneself\textquotedblright{}: to wash your own face and hands.
\end{cousinweb}

\begin{grammarlens}[title={What you've built}]
One more word for this chapter.
\end{grammarlens}

\subsection*{Your turn}
\begin{itemize}
\raggedright
  \item \emph{Say it:} \emph{umyvátsya}
  \item \emph{Say it:} \emph{umyvátsya}, once more
  \item \emph{Say it:} say \emph{umyvátsya}
  \item \emph{Recall:} say the Russian for to drive, then the Russian for to get dressed, then say \emph{umyvátsya} again
\end{itemize}

\begin{tcolorbox}[breakable,colback=teal!4,colframe=teal!35!black,title={Before you move on}]
What does \ru{умываться} mean? (\textquotedblleft{}To wash (oneself)\textquotedblright{}.) Say it once more.
\end{tcolorbox}
\section[\texorpdfstring{\ru{вставать} (vstavát)}{вставать (vstavát)}]{\ru{вставать} (vstavát) — to get up}

\label{lesson:RU-C132-to-get-up}



\begin{quote}
Before the new one: say the Russian for to get dressed, then the Russian for to wash.
\end{quote}

\subsection*{You'll want to know: \ru{вставать}}
\textbf{\ru{вставать}} (\emph{vstavát}) — \textquotedblleft{}to get up\textquotedblright{}.

\begin{cousinweb}
\textbf{\ru{вставать}} is \emph{\ru{вс}-}, \textquotedblleft{}up\textquotedblright{}, with \emph{-\ru{ставать}}, from \emph{\ru{стать}}, \textquotedblleft{}to stand\textquotedblright{}, a cousin of Latin \emph{stāre}, \textquotedblleft{}to stand\textquotedblright{} (Vasmer).
\end{cousinweb}

\begin{grammarlens}[title={What you've built}]
One more word for this chapter.
\end{grammarlens}

\subsection*{Your turn}
\begin{itemize}
\raggedright
  \item \emph{Say it:} \emph{vstavát}
  \item \emph{Say it:} \emph{vstavát}, once more
  \item \emph{Say it:} say \emph{vstavát}
  \item \emph{Recall:} say the Russian for to get dressed, then the Russian for to wash, then say \emph{vstavát} again
\end{itemize}

\begin{tcolorbox}[breakable,colback=teal!4,colframe=teal!35!black,title={Before you move on}]
What does \ru{вставать} mean? (\textquotedblleft{}To get up\textquotedblright{}.) Say it once more.
\end{tcolorbox}
\section[\texorpdfstring{\ru{садиться} (sadítsya)}{садиться (sadítsya)}]{\ru{садиться} (sadítsya) — to sit down}

\label{lesson:RU-C132-to-sit-down}



\begin{quote}
Before the new one: say the Russian for to wash, then the Russian for to get up.
\end{quote}

\subsection*{You'll want to know: \ru{садиться}}
\textbf{\ru{садиться}} (\emph{sadítsya}) — \textquotedblleft{}to sit down\textquotedblright{}.

\begin{cousinweb}
\textbf{\ru{садиться}} is \emph{\ru{садить}}, \textquotedblleft{}to seat\textquotedblright{}, with \emph{-\ru{ся}}, \textquotedblleft{}oneself\textquotedblright{}: to seat yourself. \emph{\ru{садить}} $\leftarrow$ Proto-Slavic \emph{*saditi}, the \textquotedblleft{}make sit\textquotedblright{} form of \emph{\ru{сидеть}}, \textquotedblleft{}to sit\textquotedblright{}; Vasmer relates it to Sanskrit \emph{sādáyati}, \textquotedblleft{}seats\textquotedblright{}, and Gothic \emph{satjan}, \textquotedblleft{}to set\textquotedblright{}.
\end{cousinweb}

\begin{grammarlens}[title={What you've built}]
One more word for this chapter.
\end{grammarlens}

\subsection*{Your turn}
\begin{itemize}
\raggedright
  \item \emph{Say it:} \emph{sadítsya}
  \item \emph{Say it:} \emph{sadítsya}, once more
  \item \emph{Say it:} say \emph{sadítsya}
  \item \emph{Recall:} say the Russian for to wash, then the Russian for to get up, then say \emph{sadítsya} again
\end{itemize}

\begin{tcolorbox}[breakable,colback=teal!4,colframe=teal!35!black,title={Before you move on}]
What does \ru{садиться} mean? (\textquotedblleft{}To sit down\textquotedblright{}.) Say it once more.
\end{tcolorbox}
\section[\texorpdfstring{\ru{падать} (pádat)}{падать (pádat)}]{\ru{падать} (pádat) — to fall}

\label{lesson:RU-C132-to-fall}



\begin{quote}
Before the new one: say the Russian for to get up, then the Russian for to sit down.
\end{quote}

\subsection*{You'll want to know: \ru{падать}}
\textbf{\ru{падать}} (\emph{pádat}) — \textquotedblleft{}to fall\textquotedblright{}.

\begin{cousinweb}
\textbf{\ru{падать}} $\leftarrow$ Proto-Slavic \emph{*padati} — inherited. Vasmer relates it to Sanskrit \emph{padyate}, \textquotedblleft{}falls, goes\textquotedblright{}, and Old English \emph{fetan}, \textquotedblleft{}to fall\textquotedblright{}.
\end{cousinweb}

\begin{grammarlens}[title={What you've built}]
One more word for this chapter.
\end{grammarlens}

\subsection*{Your turn}
\begin{itemize}
\raggedright
  \item \emph{Say it:} \emph{pádat}
  \item \emph{Say it:} \emph{pádat}, once more
  \item \emph{Say it:} say \emph{pádat}
  \item \emph{Recall:} say the Russian for to get up, then the Russian for to sit down, then say \emph{pádat} again
\end{itemize}

\begin{tcolorbox}[breakable,colback=teal!4,colframe=teal!35!black,title={Before you move on}]
What does \ru{падать} mean? (\textquotedblleft{}To fall\textquotedblright{}.) Say it once more.
\end{tcolorbox}
